\documentclass[11pt]{article}
\usepackage[margin=1in]{geometry}
\usepackage{amsmath,amssymb,amsthm,mathtools}
\usepackage[hidelinks]{hyperref}
\newtheorem{theorem}{Theorem}
\newtheorem{lemma}[theorem]{Lemma}
\newtheorem{corollary}[theorem]{Corollary}
\newcommand{\Q}{\mathbb Q}
\newcommand{\F}{\mathbb F}
\newcommand{\N}{\mathbb N}
\title{Natural density one squarefreeness for the exponent-six factorial moments}
\author{User-supplied asymptotic argument checked and integrated; no priority claim}
\date{16 November 2026}
\begin{document}
\maketitle
Put $F_n(T)=\sum_{j=0}^n\binom nj(6j)!T^j$. All factorials have ordinary
Gamma shape one. This note combines the retained prime-gap and shifted
Newton lemmas with credited published analytic number theory. It proves
natural density one, with no all-index squarefreeness or full SFC claim.

\begin{theorem}\label{thm:density}
Unconditionally,
\[
 \#\{1\le n\le X:F_n\text{ is not squarefree over }\Q\}=o(X).
\]
No fixed multiplier or offset cutoff occurs in this asymptotic statement.
It does not assert squarefreeness for every multiplier, offset or index.
\end{theorem}

\section{Published inputs and their exact scopes}
The pinned author-hosted journal Online First version of
\href{sources/filaseta-rocks-2016-published.pdf}{Filaseta--Rocks},
Theorem~1.1, proves
\[
 \liminf_{H\to\infty}\frac{|I\cap[1,H]|}{H}\ge\log2,
 \quad I=\{i\ge1:F_i\text{ is irreducible over }\Q\}.
\]
Their $f_{n,m}$ is exactly this $F_n$ when $m=6$.
In their Section~3, the set $B$ is defined by the existence of a prime
$n-\sqrt n<q\le n$, and equation~(3.3) gives $B$ natural density one.
These unconditional analytic results are imported, not independently
reproved from their deeper prime-distribution inputs.

The pinned publisher version of
\href{sources/matomaki-radziwill-2016-annals.pdf}{Matom\"aki--Radziwi\l\l},
Theorem~1, pp.~1015--1016, gives absolute $C,C'>1$ such that for every
multiplicative $f:\N\to[-1,1]$, $2\le h\le X$ and $\delta>0$,
\[
 \left|\frac1h\sum_{x\le u\le x+h}f(u)
        -\frac1X\sum_{X\le u\le2X}f(u)\right|
 \le\delta+C'\frac{\log\log h}{\log h}
\]
except at most
\[
 CX\left(\frac{(\log h)^{1/3}}{\delta^2h^{\delta/25}}
       +\frac1{\delta^2(\log X)^{1/50}}\right)
 \tag{1}
\]
integer starts $x\in[X,2X]$. One may take $C'=20000$.
The text immediately following the theorem explicitly states uniformity
in $h,\delta,f$. In particular the function may depend on $X$.
This deep theorem is a credited import; finite controls do not prove it.

\section{The elementary long mean and the short-interval hitting step}
For $Y=(2X)^{2/3}$ put $f_X(u)=1$ if all prime divisors of $u$ are at
most $Y$, and zero otherwise, with $f_X(1)=1$.
It is a real, bounded, completely multiplicative function.
Since $Y^2>2X$, every integer $u\le2X$ has at most one prime divisor
above $Y$, counted to the first power. Thus the exact disjoint-multiple
identity and floor errors give
\[
 \frac1X\sum_{X\le u\le2X}f_X(u)
 =1-\sum_{Y<p\le2X}\frac1p+O\left(\frac{\pi(2X)+1}{X}\right).
\]
Partial summation from the classical prime number theorem gives, uniformly
on this range with fixed logarithmic endpoint ratio,
\[
 \sum_{Y<p\le2X}\frac1p
 =\log\frac{\log(2X)}{\log Y}+o(1)=\log(3/2)+o(1).
\]
Indeed the relative PNT error is uniformly $o(1)$ for all $t\ge Y$,
and its integral is bounded by that error times
$\int_Y^{2X}dt/(t\log t)=\log(3/2)$.
Consequently the long mean tends to $1-\log(3/2)<3/5$.
This also identifies $\rho(3/2)$ without importing a separate smooth-number
theorem. Exact two-term positive logarithm series give
\[
 \log2>2\left(\frac13+\frac1{81}\right)=\frac{56}{81}>\frac23,
 \qquad
 \log(3/2)>2\left(\frac15+\frac1{375}\right)
 =\frac{152}{375}>\frac25.
\]

Fix $\delta=1/100$. For every sufficiently large fixed integer $H$,
Filaseta--Rocks gives $|I_H|>2H/3$, where $I_H=I\cap[1,H]$, and
$C'\log\log H/\log H<1/100$.
For all sufficiently large $X$ the long mean is below $3/5$.
Apply Theorem~1 to $f_X$ with $h=H$ and set $x=n-H$.
For $X+H\le n\le2X$, this start lies in the theorem's range.
The desired backward block $n-H,\ldots,n-1$ is a subset of the inclusive
forward block $x,\ldots,x+H$; nonnegativity permits dropping its last term.
Only the $O(H)$ indices $X\le n<X+H$ are discarded. Hence
\[
 \sum_{i=1}^H f_X(n-i)<\frac{31H}{50}<\frac{2H}{3} \tag{2}
\]
except at most $(\eta(H)+o_H(1))X$ indices $n\in[X,2X]$, where
\[
 \eta(H)=\frac{C(\log H)^{1/3}}{\delta^2H^{\delta/25}}
 \longrightarrow0.
\]
If every $n-i$ for $i\in I_H$ were smooth, the sum in (2) would be
at least $|I_H|>2H/3$, a contradiction. Thus some $i\in I_H$ has
\[
 p=P^+(n-i)>(2X)^{2/3}. \tag{3}
\]
The offset is selected from this fixed finite set; no independence among
individual shifts or irreducibility tests at growing offsets is assumed.

\section{Finite discriminants and the local contradiction}
For the fixed $H$, every $F_i$ with $i\in I_H$ is squarefree in
characteristic zero. Its integer discriminant $\Delta_i$ is nonzero.
There is therefore a finite $T_H$ such that for every prime
$p>T_H$, all these reductions are squarefree and preserve their degrees.
For example, $T_H=\max(6H,\max_{i\in I_H}|\Delta_i|)$ is sufficient.
This conservative existence threshold requires no uniform assertion about
all offsets or all first-subresultant integers $C_i$.

For fixed $H$ and sufficiently large $X$, (3) gives $p>T_H$, $p>6H$
and $p^2>6n$, since $n\le2X$ and $(2X)^{4/3}>12X$ eventually.
Writing $n=ap+i$, we have $a\ge1$ and $p^2>6n>6ap$,
so $p>6\max(a,i)$. The multiplier may grow without bound with $X$.

Now also require $n\in B$, discarding only $o(X)$ dyadic indices.
For the largest prime $q_n\le n$, the retained
\href{PRIME_GAPS_AND_SMALL_KERNELS.tex}{prime-gap theorem} gives, for
any rational square divisor $g^2\mid F_n$ and large $n$,
\[
 k=\deg g\le\left\lfloor\frac{n-q_n}{2}\right\rfloor
 <\frac{\sqrt n}{2}<p. \tag{4}
\]
The stronger even rounding is available, but is unnecessary here.
Choose $g$ primitive integral. By Gauss's lemma its square has an integer
quotient in $F_n$, and $F_n(0)=1$ forces $g(0)=\pm1$.
The retained \href{SLOPE_COVERAGE_AND_SHIFTED_WINDOWS.tex}{shifted polygon}
at $p$ has a horizontal edge of length $i$ followed by an edge of
reduced slope denominator $p$. Each positive sloped contribution to a
factor degree is a positive multiple of $p$. By (4), $g$ lies entirely
on the horizontal edge; its leading coefficient is a $p$-adic unit and
its reduction preserves degree $k$.

Lucas's congruence and the factorial coefficients give the full identity
\[
 \overline{F_n}=\overline{F_i}\quad\text{in }\F_p[T]. \tag{5}
\]
For $j\le i$, $\binom nj\equiv\binom ij\pmod p$.
For $i<j<p/6$, the binomial coefficient vanishes by Lucas, and for
$j\ge\lceil p/6\rceil$ the factorial $(6j)!$ vanishes modulo $p$.
Reducing the integer factorization therefore yields
$\bar g^2\mid\overline{F_i}$ with $\deg\bar g=k$.
The chosen $p>T_H$ makes the right side squarefree, so $k=0$.
Every nonconstant repeated rational irreducible factor is thereby excluded.

\section{Order of limits and ordinary natural density}
For every sufficiently large fixed $H$ the preceding argument proves
\[
 \limsup_{X\to\infty}\frac{\#\{X\le n\le2X:F_n\text{ not squarefree}\}}{X}
 \le\eta(H).
\]
All finite bad-discriminant primes and offset-dependent thresholds are
passed while $H$ remains fixed. Letting $H\to\infty$ now proves zero
dyadic exceptional density. Given $\varepsilon>0$, all sufficiently large
dyadic blocks have exceptional count at most $\varepsilon$ times their
length. Summing blocks down from $X$ and retaining the bounded initial
remainder proves $o(X)$ exceptional indices up to $X$, establishing
Theorem~\ref{thm:density}. Half-open blocks avoid endpoint overlaps.
The logical order is fixed $H$, then $X\to\infty$, then $H\to\infty$;
no uncontrolled substitution $H=H(X)$ is used.

\begin{corollary}
A set of terminal indices of natural density one simultaneously certifies
at least one nonzero moment in the pair $n-1,n$ for every nonzero complex
phase $A+BY^6$ in ordinary shape one.
\end{corollary}
\begin{proof}
Use the retained identity $F_{n-1}=F_n-(T/n)F_n'$ and $F_n(0)=1$.
Squarefreeness excludes a common root. The cases $A=0$ or $B=0$ are
immediate, and the others use $T=B/A$.
\end{proof}
A possible density-zero exceptional set remains. Full all-index SFC,
all-index residual-degree coverage and other Gamma shapes are not settled.
This is stronger than the imported lower irreducibility density bound,
while it makes no irreducibility-density-one or novelty claim.

\section{Executable evidence boundary}
\href{density_one_windows.py}{Finite local certificates} combine an
explicit shifted prime with the global degree bound and a checked modular
squarefree kernel at any supplied offset. They do not search for the
analytic offset sets implicitly.
\href{verify_density_one.py}{Exact controls} check coefficient reduction,
finite-field Bezout identities, the disjoint-large-prime counting identity,
source pins and rational constant gaps. They do not numerically establish
an asymptotic density theorem. The written argument above and the credited
published inputs supply its universal evidence. Independent replay,
formal verification and external review remain unclaimed.
\end{document}
